\documentclass[11pt,a4paper]{article}

% ==========================================================
%  MANUAL DEL PROGRAMADOR
%  StockPro - Sistema de Gestión de Inventario
% ==========================================================

\usepackage[utf8]{inputenc}
\usepackage[T1]{fontenc}
\usepackage[spanish,es-tabla]{babel}
\usepackage{lmodern}
\usepackage[margin=2.5cm]{geometry}
\usepackage{graphicx}
\usepackage{xcolor}
\usepackage{listings}
\usepackage{longtable}
\usepackage{array}
\usepackage{booktabs}
\usepackage{enumitem}
\usepackage{titlesec}
\usepackage{fancyhdr}
\usepackage{tabularx}
\usepackage[hidelinks]{hyperref}

% ---------- Colores ----------
\definecolor{azulcorp}{RGB}{0,51,102}
\definecolor{grisclaro}{RGB}{245,245,245}
\definecolor{verdecom}{RGB}{0,128,0}
\definecolor{cadenas}{RGB}{163,21,21}
\definecolor{palabraclave}{RGB}{0,0,200}

% ---------- Hipervínculos ----------
\hypersetup{
    colorlinks=true,
    linkcolor=azulcorp,
    urlcolor=azulcorp,
    citecolor=azulcorp
}

% ---------- Encabezado y pie de página ----------
\pagestyle{fancy}
\fancyhf{}
\fancyhead[L]{\small Manual del Programador}
\fancyhead[R]{\small StockPro --- Gestión de Inventario}
\fancyfoot[C]{\thepage}
\renewcommand{\headrulewidth}{0.4pt}
\renewcommand{\footrulewidth}{0.2pt}

% ---------- Estilo de títulos ----------
\titleformat{\section}
  {\color{azulcorp}\Large\bfseries}{\thesection}{1em}{}
\titleformat{\subsection}
  {\color{azulcorp}\large\bfseries}{\thesubsection}{1em}{}

% ---------- Configuración de listings (código) ----------
\lstdefinestyle{codigo}{
    backgroundcolor=\color{grisclaro},
    basicstyle=\ttfamily\footnotesize,
    keywordstyle=\color{palabraclave}\bfseries,
    commentstyle=\color{verdecom}\itshape,
    stringstyle=\color{cadenas},
    numberstyle=\tiny\color{gray},
    numbers=left,
    numbersep=8pt,
    showstringspaces=false,
    breaklines=true,
    frame=single,
    rulecolor=\color{gray!40},
    tabsize=2,
    captionpos=b
}
\lstset{style=codigo}

% Lenguajes adicionales
\lstdefinelanguage{TypeScript}{
  keywords={import, export, from, const, let, var, function, class, interface,
            extends, implements, return, if, else, for, while, new, this,
            public, private, readonly, async, await, type, signal, computed, inject},
  sensitive=true,
  comment=[l]{//},
  morecomment=[s]{/*}{*/},
  morestring=[b]',
  morestring=[b]"
}

% ==========================================================
\begin{document}

% ----------------------- PORTADA --------------------------
\begin{titlepage}
    \centering
    \vspace*{1.5cm}

    {\color{azulcorp}\rule{\linewidth}{1.2pt}}\\[0.4cm]
    {\Huge\bfseries\color{azulcorp} Manual del Programador}\\[0.3cm]
    {\LARGE StockPro \textemdash{} Sistema de Gestión de Inventario}\\[0.2cm]
    {\color{azulcorp}\rule{\linewidth}{1.2pt}}\\[1.5cm]

    \vfill

    {\large
    \begin{tabular}{rl}
        \textbf{Proyecto:} & StockPro --- Gestión de Inventario \\[4pt]
        \textbf{Versión del sistema:} & 1.0.0 \\[4pt]
        \textbf{Tipo de documento:} & Manual técnico / del programador \\[4pt]
        \textbf{Arquitectura:} & Cliente--Servidor (SPA + API REST) \\[4pt]
        \textbf{Backend:} & Java 17 / Spring Boot 3.2.5 \\[4pt]
        \textbf{Frontend:} & Angular 21 \\[4pt]
        \textbf{Base de datos:} & MySQL 8.0 \\[4pt]
    \end{tabular}
    }

    \vfill

    {\large \today}

\end{titlepage}

% --------------------- ÍNDICE -----------------------------
\tableofcontents
\newpage

% ==========================================================
\section{Introducción}
% ==========================================================

Este documento constituye el \textbf{Manual del Programador} de
\textbf{StockPro}, un sistema de gestión de inventario y punto de venta. Su
finalidad es servir de guía técnica para
cualquier desarrollador que deba \emph{instalar, comprender, mantener o
extender} el sistema. A diferencia de un manual de usuario, este manual no
describe cómo operar la aplicación, sino cómo está construida internamente:
su arquitectura, las herramientas empleadas, la estructura del código fuente,
el modelo de datos y los procedimientos de compilación y despliegue.

\subsection{Objetivo del manual}

Proporcionar la información técnica necesaria para que un programador pueda:

\begin{itemize}[noitemsep]
    \item Preparar el entorno de desarrollo y poner en marcha el sistema.
    \item Entender la arquitectura general y la organización del código.
    \item Localizar y modificar las distintas capas (controladores,
          servicios, repositorios, componentes).
    \item Comprender el modelo de datos y las reglas de negocio.
    \item Compilar, empaquetar y desplegar la aplicación.
\end{itemize}

\subsection{Alcance}

El sistema gestiona el inventario de un establecimiento comercial e incluye
módulos de productos, categorías, proveedores, clientes, impuestos, control
de lotes con fechas de vencimiento, movimientos de stock (entradas, salidas
y mermas), ventas con múltiples productos, estadísticas y administración de
usuarios con dos roles diferenciados (\texttt{ADMIN} y \texttt{VENDEDOR}).

\subsection{Descripción general del sistema}

La aplicación está dividida en tres componentes independientes que se
comunican entre sí:

\begin{itemize}[noitemsep]
    \item Un \textbf{frontend} (interfaz web SPA) desarrollado en Angular,
          servido por Nginx.
    \item Un \textbf{backend} (API REST) desarrollado en Spring Boot, que
          expone los servicios bajo la ruta \texttt{/api}.
    \item Una \textbf{base de datos} relacional MySQL que persiste la
          información.
\end{itemize}

La autenticación se realiza mediante \textbf{tokens JWT} (JSON Web Tokens) sin
estado de sesión en el servidor; el token viaja en la cabecera
\texttt{Authorization} de cada petición.

% ==========================================================
\section{Requisitos y herramientas}
% ==========================================================

\subsection{Requisitos de hardware (recomendados)}

\begin{itemize}[noitemsep]
    \item Procesador de 64 bits, doble núcleo o superior.
    \item 4 GB de memoria RAM como mínimo (8 GB recomendado para compilar).
    \item 2 GB de espacio libre en disco.
\end{itemize}

\subsection{Requisitos de software}

\begin{longtable}{>{\bfseries}p{4.5cm} p{4cm} p{5cm}}
\toprule
\textbf{Herramienta} & \textbf{Versión} & \textbf{Uso} \\
\midrule
\endhead
Java JDK & 17 & Compilación y ejecución del backend \\
Maven & 3.9+ & Gestión de dependencias del backend \\
Node.js & 20+ & Compilación del frontend \\
npm & 11+ & Gestión de paquetes del frontend \\
Angular CLI & 21+ & Desarrollo y \emph{build} del frontend \\
MySQL & 8.0 & Motor de base de datos \\
Docker / Docker Compose & Reciente & Despliegue en contenedores (opcional) \\
Git & Reciente & Control de versiones \\
\bottomrule
\end{longtable}

\subsection{Tecnologías y librerías principales}

\textbf{Backend (Spring Boot 3.2.5, Java 17):}
\begin{itemize}[noitemsep]
    \item \texttt{spring-boot-starter-web} --- API REST.
    \item \texttt{spring-boot-starter-data-jpa} --- persistencia con Hibernate/JPA.
    \item \texttt{spring-boot-starter-security} --- seguridad y autenticación.
    \item \texttt{spring-boot-starter-validation} --- validación de datos de entrada.
    \item \texttt{mysql-connector-j} --- conector JDBC para MySQL.
    \item \texttt{jjwt} (0.12.5) --- generación y validación de tokens JWT.
    \item \texttt{lombok} (1.18.36) --- reducción de código repetitivo (getters, setters, \emph{builders}).
\end{itemize}

\textbf{Frontend (Angular 21):}
\begin{itemize}[noitemsep]
    \item \texttt{@angular/core}, \texttt{@angular/router}, \texttt{@angular/forms}, \texttt{@angular/common/http}.
    \item \texttt{rxjs} --- programación reactiva.
    \item Componentes \emph{standalone} y \emph{signals} (API moderna de Angular).
    \item Nginx --- servidor web y \emph{proxy} hacia el backend.
\end{itemize}

% ==========================================================
\section{Arquitectura del sistema}
% ==========================================================

El sistema sigue una arquitectura \textbf{cliente--servidor} de tres capas
físicas. El backend implementa internamente una arquitectura en capas
(\emph{layered architecture}): Controlador $\rightarrow$ Servicio
$\rightarrow$ Repositorio $\rightarrow$ Base de datos.

\begin{center}
\footnotesize
\begin{tabular}{|c|}
\hline
\textbf{NAVEGADOR (Cliente)} \\
Angular SPA --- componentes, servicios, guards, interceptor JWT \\
\hline
$\downarrow$\quad HTTP / JSON \quad (Bearer token)\quad $\uparrow$ \\
\hline
\textbf{NGINX} (puerto 80) --- sirve la SPA y hace \emph{proxy} de \texttt{/api} \\
\hline
$\downarrow$\quad \texttt{proxy\_pass} \quad $\uparrow$ \\
\hline
\textbf{BACKEND Spring Boot} (puerto 8080) \\
Controller $\rightarrow$ Service $\rightarrow$ Repository \\
Filtro JWT + Spring Security \\
\hline
$\downarrow$\quad JDBC / JPA \quad $\uparrow$ \\
\hline
\textbf{MySQL 8.0} --- base de datos \texttt{inventario\_escolar} \\
\hline
\end{tabular}
\end{center}

\subsection{Capas del backend}

\begin{longtable}{>{\bfseries}p{3cm} p{11.5cm}}
\toprule
\textbf{Capa} & \textbf{Responsabilidad} \\
\midrule
\endhead
Controller & Expone los \emph{endpoints} REST, recibe peticiones HTTP,
            valida la entrada y delega en los servicios. Paquete
            \texttt{com.inventario.controller}. \\
Service & Contiene la lógica de negocio y las transacciones
          (\texttt{@Transactional}). Paquete \texttt{com.inventario.service}. \\
Repository & Acceso a datos mediante Spring Data JPA. Interfaces que
             extienden \texttt{JpaRepository}. Paquete
             \texttt{com.inventario.repository}. \\
Entity & Entidades JPA mapeadas a tablas. Paquete
         \texttt{com.inventario.entity}. \\
DTO & Objetos de transferencia de datos para peticiones/respuestas.
      Paquete \texttt{com.inventario.dto}. \\
Security & Configuración de seguridad, filtro y utilidades JWT. Paquetes
           \texttt{com.inventario.security} y \texttt{com.inventario.config}. \\
Exception & Excepciones personalizadas y manejador global. Paquete
            \texttt{com.inventario.exception}. \\
\bottomrule
\end{longtable}

% ==========================================================
\section{Estructura de directorios}
% ==========================================================

A continuación se muestra la organización del repositorio.

\begin{lstlisting}[language=bash,caption={Estructura raíz del proyecto}]
StockPro/
|-- backend/                 # API REST (Spring Boot)
|   |-- src/main/java/com/inventario/
|   |   |-- InventarioApplication.java   # Clase principal
|   |   |-- config/          # SecurityConfig
|   |   |-- controller/      # Controladores REST
|   |   |-- dto/             # Objetos de transferencia
|   |   |-- entity/          # Entidades JPA
|   |   |-- enums/           # Rol, TipoMovimiento
|   |   |-- exception/       # Excepciones + manejador global
|   |   |-- repository/      # Repositorios Spring Data JPA
|   |   |-- security/        # JWT y UserDetailsService
|   |   `-- service/         # Logica de negocio
|   |-- src/main/resources/application.properties
|   |-- pom.xml              # Dependencias Maven
|   `-- Dockerfile
|-- frontend/                # SPA (Angular)
|   |-- src/app/
|   |   |-- core/            # guards, interceptors, layout, models, services
|   |   |-- features/        # modulos funcionales (vistas)
|   |   |-- app.routes.ts    # Rutas
|   |   `-- app.config.ts    # Configuracion de la aplicacion
|   |-- src/environments/    # environment.ts / environment.prod.ts
|   |-- nginx.conf
|   |-- package.json
|   `-- Dockerfile
|-- docker-compose.yml       # Orquestacion de los 3 servicios
`-- INSTRUCCIONES-WINDOWS.md
\end{lstlisting}

\subsection{Estructura del frontend (detalle)}

\begin{lstlisting}[language=bash,caption={Organización de \texttt{frontend/src/app}}]
app/
|-- core/
|   |-- guards/        # auth.guard.ts, role.guard.ts
|   |-- interceptors/  # auth.interceptor.ts (inyecta el token)
|   |-- layout/        # admin-layout, vendedor-layout
|   |-- models/        # interfaces TypeScript (DTOs del cliente)
|   `-- services/      # un servicio HTTP por entidad
`-- features/
    |-- auth/          # login, register
    |-- dashboard/     # admin-dashboard, vendedor-dashboard
    |-- productos/  inventario/  movimientos/  ventas/
    |-- proveedores/ clientes/  categorias/  lotes/  usuarios/
\end{lstlisting}

% ==========================================================
\section{Configuración del entorno de desarrollo}
% ==========================================================

\subsection{Backend}

\begin{enumerate}[noitemsep]
    \item Instalar JDK 17 y Maven 3.9+.
    \item Crear la base de datos en MySQL:
\begin{lstlisting}[language=SQL]
CREATE DATABASE inventario_escolar
  CHARACTER SET utf8mb4 COLLATE utf8mb4_unicode_ci;
\end{lstlisting}
    \item Revisar la configuración en
          \texttt{backend/src/main/resources/application.properties}
          (ver sección \ref{sec:config}).
    \item Compilar y ejecutar:
\begin{lstlisting}[language=bash]
cd backend
mvn clean package          # genera target/*.jar
mvn spring-boot:run        # o: java -jar target/inventario-backend-1.0.0.jar
\end{lstlisting}
    El backend queda disponible en \texttt{http://localhost:8080}.
\end{enumerate}

\subsection{Frontend}

\begin{enumerate}[noitemsep]
    \item Instalar Node.js 20+ y Angular CLI.
    \item Instalar dependencias y arrancar el servidor de desarrollo:
\begin{lstlisting}[language=bash]
cd frontend
npm install
npm start                  # ng serve -> http://localhost:4200
\end{lstlisting}
    \item En desarrollo, el frontend apunta a \texttt{http://localhost:8080/api}
          (ver \texttt{src/environments/environment.ts}).
\end{enumerate}

\subsection{Configuración (application.properties)}\label{sec:config}

\begin{lstlisting}[language=bash,caption={Parámetros configurables del backend}]
server.port=8080
spring.datasource.url=jdbc:mysql://localhost:3306/inventario_escolar?...
spring.datasource.username=root
spring.jpa.hibernate.ddl-auto=update     # esquema autogenerado por Hibernate
spring.jpa.show-sql=true
jwt.secret=<clave secreta de al menos 256 bits>
jwt.expiration=86400000                  # 24 h en milisegundos
\end{lstlisting}

\begin{itemize}[noitemsep]
    \item \texttt{ddl-auto=update}: Hibernate crea/actualiza las tablas
          automáticamente a partir de las entidades. En el despliegue con
          Docker se usa \texttt{create-drop} (las tablas se recrean al
          arrancar y se borran al apagar; \emph{los datos no persisten} entre
          reinicios).
    \item \texttt{jwt.secret} y \texttt{jwt.expiration} controlan la firma y la
          caducidad de los tokens. Pueden sobrescribirse con variables de
          entorno (\texttt{JWT\_SECRET}, \texttt{JWT\_EXPIRATION}).
\end{itemize}

% ==========================================================
\section{Modelo de datos}
% ==========================================================

El esquema relacional se genera automáticamente a partir de las entidades JPA.
La siguiente tabla resume las entidades y sus tablas asociadas.

\begin{longtable}{>{\bfseries}p{3cm} p{3.5cm} p{8cm}}
\toprule
\textbf{Entidad} & \textbf{Tabla} & \textbf{Descripción} \\
\midrule
\endhead
Usuario & \texttt{usuarios} & Usuarios del sistema con rol \texttt{ADMIN} o \texttt{VENDEDOR}. \\
Categoria & \texttt{categorias} & Clasificación de productos. \\
Proveedor & \texttt{proveedores} & Proveedores de productos. \\
Impuesto & \texttt{impuestos} & Impuestos aplicables (nombre y porcentaje). \\
Producto & \texttt{productos} & Productos del inventario (SKU único, stock, precios). \\
Lote & \texttt{lotes} & Lotes de un producto con fecha de vencimiento. \\
Movimiento & \texttt{movimientos} & Entradas, salidas y mermas de stock. \\
Cliente & \texttt{clientes} & Clientes a quienes se realizan ventas. \\
Venta & \texttt{ventas} & Cabecera de venta (totales, cliente, usuario). \\
DetalleVenta & \texttt{detalle\_venta} & Líneas de cada venta (producto, cantidad, subtotal). \\
\bottomrule
\end{longtable}

\subsection{Relaciones principales}

\begin{itemize}[noitemsep]
    \item \texttt{Producto} $\to$ \texttt{Categoria}, \texttt{Impuesto},
          \texttt{Proveedor} (\texttt{ManyToOne}).
    \item \texttt{Lote} $\to$ \texttt{Producto} (\texttt{ManyToOne}).
    \item \texttt{Movimiento} $\to$ \texttt{Producto}, \texttt{Usuario}
          (\texttt{ManyToOne}).
    \item \texttt{Venta} $\to$ \texttt{Cliente}, \texttt{Usuario}
          (\texttt{ManyToOne}); \texttt{Venta} $\to$ \texttt{DetalleVenta}
          (\texttt{OneToMany}, \texttt{cascade=ALL}, \texttt{orphanRemoval}).
    \item \texttt{DetalleVenta} $\to$ \texttt{Producto} (\texttt{ManyToOne}).
\end{itemize}

\subsection{Diccionario de datos (entidades clave)}

\textbf{Producto} (\texttt{productos})
\begin{longtable}{>{\ttfamily}p{3.2cm} p{3.5cm} p{7cm}}
\toprule
\textbf{\normalfont Campo} & \textbf{\normalfont Tipo} & \textbf{\normalfont Descripción} \\
\midrule
\endhead
id & Integer (PK) & Identificador autogenerado. \\
nombre & String(150) & Nombre del producto. Obligatorio. \\
descripcion & TEXT & Descripción larga. \\
sku & String(50) & Código único de producto. Obligatorio y único. \\
codigoBarras & String(100) & Código de barras (se sincroniza con el lote). \\
categoria\_id & FK & Referencia a \texttt{categorias}. \\
precioCompra & Decimal(10,2) & Precio de compra. Obligatorio. \\
precioVenta & Decimal(10,2) & Precio de venta. Obligatorio. \\
stock & Integer & Existencias actuales (por defecto 0). \\
stockMinimo & Integer & Umbral de stock bajo (por defecto 0). \\
ubicacion & String(100) & Ubicación física. \\
impuesto\_id & FK & Referencia a \texttt{impuestos}. \\
proveedor\_id & FK & Referencia a \texttt{proveedores}. \\
\bottomrule
\end{longtable}

\textbf{Usuario} (\texttt{usuarios})
\begin{longtable}{>{\ttfamily}p{3.2cm} p{3.5cm} p{7cm}}
\toprule
\textbf{\normalfont Campo} & \textbf{\normalfont Tipo} & \textbf{\normalfont Descripción} \\
\midrule
\endhead
id & Integer (PK) & Identificador. \\
nombre & String(100) & Nombre del usuario. \\
correo & String(100) & Correo único, usado como \emph{login}. \\
password & String & Contraseña cifrada con BCrypt (no se serializa). \\
rol & Enum & \texttt{ADMIN} o \texttt{VENDEDOR}. \\
fechaCreacion & DateTime & Fecha de alta (autogenerada). \\
\bottomrule
\end{longtable}

\textbf{Venta} (\texttt{ventas}) y \textbf{DetalleVenta} (\texttt{detalle\_venta})
\begin{longtable}{>{\ttfamily}p{3.2cm} p{3.5cm} p{7cm}}
\toprule
\textbf{\normalfont Campo} & \textbf{\normalfont Tipo} & \textbf{\normalfont Descripción} \\
\midrule
\endhead
id & Integer (PK) & Identificador de venta. \\
cliente\_id & FK & Cliente (opcional). \\
usuario\_id & FK & Vendedor que registra la venta. \\
fecha & DateTime & Fecha autogenerada. \\
subtotal & Decimal(10,2) & Suma de líneas sin impuesto. \\
impuesto & Decimal(10,2) & Impuesto total. \\
total & Decimal(10,2) & Subtotal + impuesto. \\
\midrule
\multicolumn{3}{l}{\textit{DetalleVenta}} \\
\midrule
producto\_id & FK & Producto vendido. \\
cantidad & Integer & Unidades vendidas. \\
precioUnitario & Decimal(10,2) & Precio de venta unitario. \\
subtotal & Decimal(10,2) & Subtotal de la línea (con impuesto). \\
\bottomrule
\end{longtable}

\textbf{Lote} (\texttt{lotes}) y \textbf{Movimiento} (\texttt{movimientos})
\begin{longtable}{>{\ttfamily}p{3.2cm} p{3.5cm} p{7cm}}
\toprule
\textbf{\normalfont Campo} & \textbf{\normalfont Tipo} & \textbf{\normalfont Descripción} \\
\midrule
\endhead
\multicolumn{3}{l}{\textit{Lote}} \\
\midrule
producto\_id & FK & Producto al que pertenece el lote. \\
numeroLote & String(50) & Identificador del lote. \\
fechaVencimiento & Date & Fecha de caducidad. \\
cantidad & Integer & Unidades del lote (0 = vendido/retirado). \\
\midrule
\multicolumn{3}{l}{\textit{Movimiento}} \\
\midrule
producto\_id & FK & Producto afectado. \\
usuario\_id & FK & Usuario responsable. \\
tipo & Enum & \texttt{ENTRADA}, \texttt{SALIDA} o \texttt{MERMA}. \\
cantidad & Integer & Unidades del movimiento. \\
motivo & String(200) & Descripción (p. ej. ``Venta \#5''). \\
fecha & DateTime & Fecha autogenerada. \\
\bottomrule
\end{longtable}

% ==========================================================
\section{Seguridad y autenticación}
% ==========================================================

La seguridad se basa en \textbf{Spring Security} con autenticación por
\textbf{JWT} y sesiones \emph{stateless}. Los componentes implicados son:

\begin{itemize}[noitemsep]
    \item \texttt{SecurityConfig} --- define la cadena de filtros, las reglas
          de autorización y la configuración CORS.
    \item \texttt{JwtUtil} --- genera y valida los tokens (firma HMAC).
    \item \texttt{JwtAuthenticationFilter} --- intercepta cada petición, extrae
          el token de la cabecera \texttt{Authorization: Bearer ...} y
          establece la autenticación en el contexto de seguridad.
    \item \texttt{CustomUserDetailsService} --- carga el usuario desde la base
          de datos por correo y le asigna la autoridad
          \texttt{ROLE\_<rol>}.
    \item \texttt{AuthService} --- gestiona el \emph{login} y el registro.
\end{itemize}

\subsection{Reglas de autorización}

\begin{lstlisting}[language=Java,caption={Reglas definidas en SecurityConfig}]
http
  .csrf(disable)
  .cors(...)                       // origenes permitidos: cualquiera (patrones)
  .sessionManagement(STATELESS)
  .authorizeHttpRequests(auth -> auth
      .requestMatchers("/api/auth/**").permitAll()      // login y registro
      .requestMatchers(OPTIONS, "/**").permitAll()       // preflight CORS
      .requestMatchers("/api/usuarios/**").hasRole("ADMIN") // solo admin
      .anyRequest().authenticated()                      // el resto, autenticado
  )
  .exceptionHandling(401)          // token invalido/expirado -> 401
  .addFilterBefore(jwtAuthFilter, UsernamePasswordAuthenticationFilter.class);
\end{lstlisting}

\subsection{Flujo de autenticación}

\begin{enumerate}[noitemsep]
    \item El cliente envía \texttt{POST /api/auth/login} con correo y contraseña.
    \item \texttt{AuthService} valida las credenciales con el
          \texttt{AuthenticationManager} y, si son correctas, genera un JWT
          con el correo (\emph{subject}) y el rol como \emph{claim}.
    \item El cliente almacena el token en \texttt{localStorage} y lo envía en
          cada petición posterior.
    \item En el servidor, \texttt{JwtAuthenticationFilter} valida el token en
          cada llamada; si es válido, autoriza la petición según el rol.
\end{enumerate}

\paragraph{Registro inicial (bootstrap).} El \emph{endpoint} público de
registro solo permite crear un usuario \texttt{ADMIN} cuando todavía no existe
ningún usuario en el sistema. A partir de ahí, los nuevos administradores se
crean desde \texttt{/api/usuarios} (protegido por rol \texttt{ADMIN}). Las
contraseñas se almacenan cifradas con \textbf{BCrypt}.

% ==========================================================
\section{API REST --- Endpoints}
% ==========================================================

Todas las rutas cuelgan de \texttt{/api} y, salvo \texttt{/api/auth/**},
requieren un token JWT válido. El formato de intercambio es JSON.

\subsection{Autenticación}
\begin{longtable}{p{2cm} p{4.5cm} p{7.5cm}}
\toprule
\textbf{Método} & \textbf{Ruta} & \textbf{Descripción} \\
\midrule
\endhead
POST & \texttt{/api/auth/login} & Inicia sesión y devuelve el token JWT. \\
POST & \texttt{/api/auth/register} & Registra un usuario (bootstrap inicial). \\
\bottomrule
\end{longtable}

\subsection{Productos}
\begin{longtable}{p{2cm} p{6.5cm} p{5.5cm}}
\toprule
\textbf{Método} & \textbf{Ruta} & \textbf{Descripción} \\
\midrule
\endhead
GET & \texttt{/api/productos} & Lista todos los productos. \\
GET & \texttt{/api/productos/\{id\}} & Obtiene un producto. \\
GET & \texttt{/api/productos/sku/\{sku\}} & Busca por SKU. \\
GET & \texttt{/api/productos/categoria/\{id\}} & Productos de una categoría. \\
GET & \texttt{/api/productos/proveedor/\{id\}} & Productos de un proveedor. \\
GET & \texttt{/api/productos/stock-bajo} & Productos con stock $\le$ mínimo. \\
GET & \texttt{/api/productos/buscar?nombre=} & Búsqueda por nombre. \\
POST & \texttt{/api/productos} & Crea un producto. \\
PUT & \texttt{/api/productos/\{id\}} & Actualiza un producto. \\
DELETE & \texttt{/api/productos/\{id\}} & Elimina un producto. \\
PATCH & \texttt{/api/productos/\{id\}/stock?cantidad=} & Ajusta el stock. \\
\bottomrule
\end{longtable}

\subsection{Ventas}
\begin{longtable}{p{2cm} p{6cm} p{6cm}}
\toprule
\textbf{Método} & \textbf{Ruta} & \textbf{Descripción} \\
\midrule
\endhead
GET & \texttt{/api/ventas} & Lista las ventas (más recientes primero). \\
GET & \texttt{/api/ventas/\{id\}} & Obtiene una venta. \\
GET & \texttt{/api/ventas/cliente/\{id\}} & Ventas de un cliente. \\
GET & \texttt{/api/ventas/usuario/\{id\}} & Ventas de un vendedor. \\
GET & \texttt{/api/ventas/fecha?inicio=\&fin=} & Ventas en un rango de fechas. \\
POST & \texttt{/api/ventas} & Registra una venta multiproducto. \\
DELETE & \texttt{/api/ventas/\{id\}} & Anula una venta (revierte stock y lotes). \\
\bottomrule
\end{longtable}

\subsection{Lotes}
\begin{longtable}{p{2cm} p{6.5cm} p{5.5cm}}
\toprule
\textbf{Método} & \textbf{Ruta} & \textbf{Descripción} \\
\midrule
\endhead
GET & \texttt{/api/lotes} & Lista todos los lotes. \\
GET & \texttt{/api/lotes/\{id\}} & Obtiene un lote. \\
GET & \texttt{/api/lotes/producto/\{id\}} & Lotes de un producto. \\
GET & \texttt{/api/lotes/vencidos} & Lotes ya vencidos. \\
GET & \texttt{/api/lotes/proximos-vencer?dias=30} & Lotes próximos a vencer. \\
GET & \texttt{/api/lotes/paginado?pagina=\&porPagina=} & Lotes paginados con estado. \\
GET & \texttt{/api/lotes/estadisticas} & Conteos (con lote, por vencer, vencidos). \\
POST & \texttt{/api/lotes} & Crea un lote (genera entrada de stock). \\
PUT & \texttt{/api/lotes/\{id\}} & Actualiza un lote. \\
DELETE & \texttt{/api/lotes/\{id\}} & Elimina un lote. \\
POST & \texttt{/api/lotes/\{id\}/retirar} & Retira por vencimiento (merma). \\
\bottomrule
\end{longtable}

\subsection{Movimientos}
\begin{longtable}{p{2cm} p{6cm} p{6cm}}
\toprule
\textbf{Método} & \textbf{Ruta} & \textbf{Descripción} \\
\midrule
\endhead
GET & \texttt{/api/movimientos} & Lista los movimientos (recientes primero). \\
GET & \texttt{/api/movimientos/\{id\}} & Obtiene un movimiento. \\
GET & \texttt{/api/movimientos/producto/\{id\}} & Movimientos de un producto. \\
GET & \texttt{/api/movimientos/tipo/\{tipo\}} & Filtra por tipo. \\
GET & \texttt{/api/movimientos/usuario/\{id\}} & Movimientos de un usuario. \\
GET & \texttt{/api/movimientos/fecha?inicio=\&fin=} & Filtra por rango de fechas. \\
POST & \texttt{/api/movimientos} & Crea un movimiento y ajusta el stock. \\
DELETE & \texttt{/api/movimientos/\{id\}} & Elimina y revierte el movimiento. \\
\bottomrule
\end{longtable}

\subsection{Otros recursos (CRUD estándar)}
Los siguientes recursos exponen el conjunto habitual de operaciones
\texttt{GET / GET\{id\} / POST / PUT / DELETE}:

\begin{longtable}{p{4cm} p{10cm}}
\toprule
\textbf{Recurso} & \textbf{Ruta base} \\
\midrule
\endhead
Categorías & \texttt{/api/categorias} \\
Proveedores & \texttt{/api/proveedores} \\
Clientes & \texttt{/api/clientes} (+ \texttt{/buscar?nombre=}) \\
Impuestos & \texttt{/api/impuestos} \\
Usuarios & \texttt{/api/usuarios} (solo \texttt{ADMIN}) \\
Estadísticas & \texttt{/api/estadisticas} (solo \texttt{GET}) \\
\bottomrule
\end{longtable}

\subsection{Formato de respuestas y errores}

Las operaciones de borrado y similares devuelven un objeto \texttt{ApiResponse}
con los campos \texttt{success} y \texttt{message}. Los errores son gestionados
de forma centralizada por \texttt{GlobalExceptionHandler}:

\begin{longtable}{>{\ttfamily}p{6cm} p{2.5cm} p{5cm}}
\toprule
\textbf{\normalfont Excepción} & \textbf{\normalfont Código} & \textbf{\normalfont Significado} \\
\midrule
\endhead
ResourceNotFoundException & 404 & Recurso no encontrado. \\
BadRequestException & 400 & Petición inválida / regla de negocio. \\
BadCredentialsException & 401 & Credenciales incorrectas. \\
MethodArgumentNotValid & 400 & Errores de validación (mapa campo--mensaje). \\
Exception (genérica) & 500 & Error interno del servidor. \\
\bottomrule
\end{longtable}

% ==========================================================
\section{Lógica de negocio (servicios clave)}
% ==========================================================

\subsection{Registro de una venta (\texttt{VentaService.create})}

Es la operación más compleja del sistema y está marcada como
\texttt{@Transactional}. Resumen del algoritmo:

\begin{enumerate}[noitemsep]
    \item Se obtiene el cliente (opcional) y el usuario (vendedor).
    \item Por cada línea de la venta (\texttt{DetalleVentaRequest}):
    \begin{itemize}[noitemsep]
        \item Se localiza el producto.
        \item Si se indica un \texttt{numeroLote}, se valida que el lote
              exista, pertenezca al producto, no esté vendido (cantidad 0) y
              tenga stock suficiente; luego se descuenta del lote.
        \item Se verifica el stock general del producto.
        \item Se calcula el subtotal y el impuesto de la línea (según el
              porcentaje del impuesto del producto, redondeo
              \texttt{HALF\_UP}).
        \item Se descuenta el stock del producto.
    \end{itemize}
    \item Se calculan los totales (subtotal, impuesto, total) y se guarda la
          venta.
    \item Por cada línea se registra un \texttt{Movimiento} de tipo
          \texttt{SALIDA} con el motivo \texttt{"Venta \#<id>"} (o
          \texttt{"Venta \#<id> - Lote: <lote>"}).
\end{enumerate}

\paragraph{Anulación de venta (\texttt{delete}).} Restaura el stock de cada
producto, devuelve la cantidad a los lotes vendidos y elimina los movimientos
asociados. Para evitar colisiones de prefijo (p. ej. ``Venta \#1'' contra
``Venta \#10''), se valida el identificador exacto del motivo.

\subsection{Movimientos de stock (\texttt{MovimientoService})}

Cada movimiento ajusta el stock del producto de forma transaccional:
\texttt{ENTRADA} lo incrementa; \texttt{SALIDA} y \texttt{MERMA} lo
decrementan (validando que haya stock suficiente). El borrado de un movimiento
revierte el efecto correspondiente.

\subsection{Gestión de lotes (\texttt{LoteService})}

\begin{itemize}[noitemsep]
    \item \textbf{Crear lote:} verifica el producto, sincroniza su código de
          barras con el número de lote y registra un movimiento de
          \texttt{ENTRADA} con motivo ``Compra'', que aumenta el stock.
    \item \textbf{Retirar por vencimiento:} registra un movimiento de
          \texttt{MERMA} por la cantidad restante del lote y fija su cantidad a
          0 (no se elimina, queda marcado).
    \item \textbf{Estado calculado:} al paginar, cada lote recibe un estado
          (\texttt{Vendido}, \texttt{Vencido}, \texttt{Por Vencer},
          \texttt{Vigente}) y los días restantes hasta el vencimiento
          (umbral de 90 días para ``Por Vencer'').
\end{itemize}

\subsection{Estadísticas (\texttt{EstadisticasService})}

Calcula los indicadores del \emph{dashboard}: total de productos, unidades
(suma de cantidades de lotes), productos con stock bajo, valor del inventario,
ventas del mes, ventas semanales, distribución de stock por categoría y
últimos movimientos.

% ==========================================================
\section{Frontend (Angular)}
% ==========================================================

El frontend es una \emph{Single Page Application} construida con Angular 21
usando componentes \emph{standalone}, \emph{lazy loading} de rutas y
\emph{signals} para el manejo de estado.

\subsection{Arranque y configuración}

La aplicación se inicia en \texttt{main.ts} mediante
\texttt{bootstrapApplication}. La configuración (\texttt{app.config.ts})
registra el \emph{router} y el \texttt{HttpClient} con el interceptor de
autenticación.

\begin{lstlisting}[language=TypeScript,caption={app.config.ts}]
export const appConfig: ApplicationConfig = {
  providers: [
    provideBrowserGlobalErrorListeners(),
    provideRouter(routes),
    provideHttpClient(withInterceptors([authInterceptor]))
  ]
};
\end{lstlisting}

\subsection{Rutas y protección por rol}

Las rutas se cargan de forma perezosa y se protegen con dos \emph{guards}:
\texttt{authGuard} (exige sesión iniciada) y \texttt{roleGuard([...])}
(exige un rol concreto). Existen dos áreas: \texttt{/admin} (rol
\texttt{ADMIN}, acceso completo) y \texttt{/vendedor} (rol \texttt{VENDEDOR},
ventas, productos, lotes y clientes).

\begin{lstlisting}[language=TypeScript,caption={Fragmento de app.routes.ts}]
{
  path: 'admin',
  loadComponent: () => import('./core/layout/admin-layout.component')
                        .then(m => m.AdminLayoutComponent),
  canActivate: [authGuard, roleGuard(['ADMIN'])],
  children: [ /* dashboard, productos, ventas, lotes, usuarios, ... */ ]
}
\end{lstlisting}

\subsection{Servicio de autenticación (\texttt{AuthService})}

Centraliza el inicio/cierre de sesión y el estado del usuario mediante
\emph{signals} (\texttt{user}, \texttt{isLoggedIn}, \texttt{rol}). El token y
los datos del usuario se persisten en \texttt{localStorage}. El método
\texttt{redirectByRole()} envía a cada usuario a su área según el rol.

\subsection{Interceptor HTTP}

\texttt{authInterceptor} añade automáticamente la cabecera
\texttt{Authorization: Bearer <token>} a cada petición y, si el backend
responde \texttt{401} (token expirado/ inválido) en una llamada que no sea de
autenticación, cierra la sesión y redirige al \emph{login}.

\begin{lstlisting}[language=TypeScript,caption={auth.interceptor.ts (esencia)}]
export const authInterceptor: HttpInterceptorFn = (req, next) => {
  const auth = inject(AuthService);
  const token = auth.getAuthorizationHeader();
  if (token) req = req.clone({ setHeaders: { Authorization: token } });
  return next(req).pipe(
    catchError((err: HttpErrorResponse) => {
      if (err.status === 401 && !req.url.includes('/auth/') && auth.isLoggedIn())
        auth.logout();
      return throwError(() => err);
    })
  );
};
\end{lstlisting}

\subsection{Servicios de dominio}

Por cada entidad existe un servicio en \texttt{core/services} que encapsula las
llamadas HTTP al backend usando \texttt{HttpClient} y las interfaces de
\texttt{core/models}. Todos construyen la URL a partir de
\texttt{environment.apiUrl}.

\begin{lstlisting}[language=TypeScript,caption={Patrón de un servicio (ProductoService)}]
@Injectable({ providedIn: 'root' })
export class ProductoService {
  private http = inject(HttpClient);
  private apiUrl = `${environment.apiUrl}/productos`;

  getAll(): Observable<Producto[]> { return this.http.get<Producto[]>(this.apiUrl); }
  create(p: Producto) { return this.http.post<Producto>(this.apiUrl, p); }
  update(id: number, p: Producto) { return this.http.put<Producto>(`${this.apiUrl}/${id}`, p); }
  delete(id: number) { return this.http.delete(`${this.apiUrl}/${id}`); }
  // ...
}
\end{lstlisting}

\subsection{Configuración por entorno}

\begin{itemize}[noitemsep]
    \item \texttt{environment.ts} (desarrollo): \texttt{apiUrl =
          'http://localhost:8080/api'}.
    \item \texttt{environment.prod.ts} (producción): \texttt{apiUrl = '/api'}
          (Nginx hace de \emph{proxy} hacia el backend).
\end{itemize}

% ==========================================================
\section{Compilación y despliegue}
% ==========================================================

\subsection{Despliegue con Docker Compose (recomendado)}

El archivo \texttt{docker-compose.yml} orquesta los tres servicios. Para
levantar todo el sistema:

\begin{lstlisting}[language=bash]
docker compose up --build
\end{lstlisting}

\begin{longtable}{>{\bfseries}p{2.8cm} p{3.2cm} p{8cm}}
\toprule
\textbf{Servicio} & \textbf{Puerto (host)} & \textbf{Descripción} \\
\midrule
\endhead
mysql & 3000 $\to$ 3306 & Base de datos MySQL 8.0 con volumen persistente. \\
backend & 6000 $\to$ 8080 & API Spring Boot. Espera a que MySQL esté sano. \\
frontend & 2200 $\to$ 80 & SPA Angular servida por Nginx. \\
\bottomrule
\end{longtable}

\paragraph{Importante.} En el despliegue Docker el backend usa
\texttt{SPRING\_JPA\_HIBERNATE\_DDL\_AUTO=create-drop}: las tablas se crean al
arrancar y se eliminan al apagar el contenedor, por lo que \emph{los datos no
se conservan} entre reinicios. Para producción real conviene cambiarlo a
\texttt{update} y montar un volumen persistente.

\subsection{Construcción manual de imágenes}

\textbf{Backend} (\emph{multi-stage} con Maven + JRE):
\begin{lstlisting}[language=bash]
# Etapa build: maven:3.9-eclipse-temurin-17 -> mvn clean package -DskipTests
# Etapa runtime: eclipse-temurin:17-jre -> java -jar app.jar
docker build -t inventario-backend ./backend
\end{lstlisting}

\textbf{Frontend} (\emph{multi-stage} con Node + Nginx):
\begin{lstlisting}[language=bash]
# Etapa build: node:20-alpine -> npm ci && npm run build
# Etapa runtime: nginx:alpine sirve dist/frontend/browser
docker build -t inventario-frontend ./frontend
\end{lstlisting}

\subsection{Configuración de Nginx}

Nginx sirve la SPA y redirige las llamadas \texttt{/api} al backend:

\begin{lstlisting}[language=bash,caption={Fragmento de nginx.conf}]
location / { try_files $uri $uri/ /index.html; }   # enrutamiento SPA
location /api { proxy_pass http://backend:8080; }  # proxy al backend
\end{lstlisting}

% ==========================================================
\section{Guía de mantenimiento y extensión}
% ==========================================================

\subsection{Cómo añadir una nueva entidad / módulo (backend)}

\begin{enumerate}[noitemsep]
    \item Crear la entidad JPA en \texttt{entity/} con sus anotaciones
          (\texttt{@Entity}, \texttt{@Table}, columnas y relaciones).
    \item Crear el repositorio en \texttt{repository/} extendiendo
          \texttt{JpaRepository<Entidad, Integer>}.
    \item Crear el servicio en \texttt{service/} con la lógica de negocio
          (anotar con \texttt{@Transactional} las operaciones de escritura).
    \item Crear el controlador en \texttt{controller/} exponiendo los
          \emph{endpoints} REST bajo \texttt{/api/<recurso>}.
    \item Si procede, añadir DTOs en \texttt{dto/} y validaciones con
          \emph{Bean Validation}.
    \item Ajustar las reglas de autorización en \texttt{SecurityConfig} si el
          recurso requiere un rol específico.
\end{enumerate}

\subsection{Cómo añadir una nueva vista (frontend)}

\begin{enumerate}[noitemsep]
    \item Crear el componente \emph{standalone} en \texttt{features/<modulo>}.
    \item Definir su interfaz (modelo) en \texttt{core/models} y su servicio
          HTTP en \texttt{core/services}.
    \item Registrar la ruta con \texttt{loadComponent} en
          \texttt{app.routes.ts}, protegiéndola con \texttt{authGuard} y
          \texttt{roleGuard} según corresponda.
\end{enumerate}

\subsection{Buenas prácticas y consideraciones}

\begin{itemize}[noitemsep]
    \item \textbf{Transacciones:} toda operación que modifique stock, lotes o
          ventas debe ser transaccional para mantener la consistencia.
    \item \textbf{Seguridad:} la clave \texttt{jwt.secret} y las credenciales
          de base de datos no deben publicarse; usar variables de entorno en
          producción.
    \item \textbf{CORS:} actualmente se aceptan todos los orígenes mediante
          patrones para facilitar el acceso desde cualquier IP; en producción
          conviene restringirlo a los dominios autorizados.
    \item \textbf{Persistencia:} revisar \texttt{ddl-auto} antes de desplegar
          en un entorno con datos reales (evitar \texttt{create-drop}).
\end{itemize}

% ==========================================================
\section{Glosario}
% ==========================================================

\begin{longtable}{>{\bfseries}p{3.5cm} p{10.5cm}}
\toprule
\textbf{Término} & \textbf{Definición} \\
\midrule
\endhead
API REST & Interfaz de programación basada en HTTP y recursos. \\
SPA & \emph{Single Page Application}; aplicación web de una sola página. \\
JWT & \emph{JSON Web Token}; token firmado para autenticación sin estado. \\
JPA / Hibernate & Especificación / implementación de mapeo objeto--relacional. \\
DTO & \emph{Data Transfer Object}; objeto para transferir datos entre capas. \\
SKU & \emph{Stock Keeping Unit}; código único de identificación de producto. \\
Stock & Cantidad de unidades disponibles de un producto. \\
Merma & Pérdida de inventario (vencimiento, daño, etc.). \\
Lote & Conjunto de unidades de un producto con una fecha de vencimiento. \\
BCrypt & Algoritmo de cifrado (hash) de contraseñas. \\
CORS & Mecanismo de control de acceso entre orígenes distintos. \\
\bottomrule
\end{longtable}

\vfill
\begin{center}
\rule{0.5\linewidth}{0.4pt}\\[0.3cm]
\small Fin del Manual del Programador --- StockPro v1.0.0
\end{center}

\end{document}
